\section{Exact exponent summary}\label{sec:ledger}

Put
\[
\begin{gathered}
S_0=\epsilon_0=k_0=\frac{30869}{5000000},\quad
l_0=2S_0,\quad j_0=p_0=3S_0,\quad \xi=\frac1{10000},\\
a_0=\frac6{125},\quad
b_0=\frac{4166537}{30000000},\quad
\Delta_{\rm exp}=3S_0.
\end{gathered}\AutoTag{eq:old:8_1}
\]
All entries below are limiting savings from \(2/3\) at this fixed
rational parameter vector. Sections~\ref{sec:long} and
\ref{sec:epsilon} handle the finite source collars and the arbitrary
analytic losses. The row estimates in Sections~\ref{sec:outer}--\ref{sec:long}
give the exact values in the second column. The symbolic counting exponents
and their evaluations are displayed in Sections~\ref{sec:joint}--\ref{sec:intermediate}
(for \(3u\), \(3t\), \(4H\), and \(4P\)). Appendix~\ref{app:arithmetic}
collects the parameter vector, exact savings, and margins above \(S_0\).

\setlength{\LTleft}{0pt}
\setlength{\LTright}{0pt}
\renewcommand{\arraystretch}{1.35}
\footnotesize
\unnumtable
The table expands the savings used in the preceding row estimates. The
minimum row and its margin are read from the exact fractions. The rounded
decimals are annotations only.
\begin{longtable}{@{}p{0.30\textwidth}p{0.53\textwidth}p{0.12\textwidth}@{}}
\toprule
Contribution & Limiting saving & Decimal\\
\midrule
\endfirsthead
\toprule
Contribution & Limiting saving & Decimal\\
\midrule
\endhead
\(O_R\) & \(30869/5000000\) & --\\
\(O_M\) & \(30869/5000000\) & --\\
\(O_{\operatorname{rad}M}\) & \(30869/5000000\) & --\\
\(Q0\) & \(2/3\) & --\\
\(Q1\) & \(30869/5000000\) & --\\
\(Q2t_{\rm int}\) & \(972179/7500000\) & --\\
\(Q2t_{\rm harm}\) & \(30869/5000000\) & --\\
\(Q2u_{\rm int}\) & \(2222179/15000000\) & --\\
\(Q2u_{\rm harm}\) & \(30869/1250000\) & --\\
\(3u_{\rm sqrt}\) & \(5102312049347/600000000000000\) & --\\
\(3u_{\rm int}\) & \(23206592049347/300000000000000\) & --\\
\(3u_{\rm core}\) & \(2473/12000\) & --\\
\(3u_{\rm small,1}\) & \(1619393/6000000\) & --\\
\(3u_{\rm small,2}\) & \(8003/500000\) & --\\
\(3t_{\rm sqrt}\) & \(388751/60000000\) & --\\
\(3t_{\rm int}\) & \(2199179/30000000\) & --\\
\(3t_{\rm axes}\) & \(2726537/15000000\) & --\\
\(4E\) & \(30869/5000000\) & --\\
\(4H\) & \(140441/5000000\) & --\\
\(4P_{\rm graph}\) & \(30869/5000000\) & --\\
\(4P_{\rm HD}\) & \(107393/15000000\) & --\\
\(4P_{\rm cover}\) & \(2/75\) & --\\
\(4P_{\rm axes}\) & \(4166537/30000000\) & --\\
\(5W_{\rm HD}\) & \(\dfrac{29390272487700347267}{4760400458585535000000}\) & \(0.006174\)\\[4pt]
\(5W_{\rm int}\) & \(\dfrac{235120131355662893}{9520800917171070000}\) & \(0.024695\)\\[4pt]
\(L_{5W}\) & \(1203113/90000000\) & \(0.013368\)\\
\bottomrule
\end{longtable}
\renewcommand{\arraystretch}{1}
\normalsize

Some entries add up within one row estimate. These are the five \(3u\)
entries, the three \(3t\) entries, the four \(4P\) entries, the two
\(Q2\) terms per chart and the three \(5W\) entries. In each row the
largest exponent controls the sum. Equivalently, the smallest saving does.
In each \(5W\) source box, Section~\ref{sec:long} selects the best chart
and assigned-corner route. The three displayed terms are the geometric,
interval and high-degree contributions that follow.
In the \(Q2\) entries, \(Q2t_{\rm int},Q2t_{\rm harm}\) are the first and
second terms of the \(Q2t\) estimate \eqref{eq:old:5_3}. Likewise,
\(Q2u_{\rm int},Q2u_{\rm harm}\) are the terms of the \(Q2u\) estimate.
Each pair is summed within its own first-match
population.
The table has 26 scalar terms. After grouping the three \(5W\)
contributions, 24 effective terms remain.

The minimum row saving is \(S_0\). It is attained by
\(O_R,O_M,O_{\operatorname{rad}M},Q1,Q2t_{\rm harm},4E\), and
\(4P_{\rm graph}\).
In the final row the smallest of the three savings is \(5W_{\rm HD}\), and
\[
5W_{\rm HD}-S_0=
\frac{128034121242821}{1190100114646383750000}
\quad\text{(about }1.076\times10^{-7}\text{)}>0.
\AutoTag{eq:old:8_2}
\]
The 72-branch minimum used for the geometric cap is
\[
L_{5W}=\frac{1203113}{90000000},\qquad
L_{5W}-S_0=\frac{647471}{90000000}>0.
\AutoTag{eq:old:8_3}
\]
Section~\ref{sec:long} proves the branch construction and the
comparison needed for \(5W\). No LP value is used as a count by itself.

\begin{proposition}[Reserve]\label{prop:reserve}
Put
\[
g_0=S_0-\frac1{165}=\frac{18677}{165000000},\qquad
\beta_0=\frac{g_0}{16}=\frac{18677}{2640000000}.
\]
Then
\[
S_0-\beta_0-\frac{g_0}{4}
=\frac1{165}+\frac{11g_0}{16}
=\frac{16205447}{2640000000}>\frac1{165}.
\AutoTag{eq:old:8_4}
\]
\end{proposition}

\begin{proof}
Put the displayed fractions over the common denominator
\(2640000000\). The loss \(\beta_0\) is charged once to each complete
analytic row estimate. In the last row, the \(g_0/4\) term covers the maximum of the
geometric, interval, and high-degree corrections.
\end{proof}
